\chapter{Criteria and the main theorem}
\label{chap:main}

Throughout, $a = r/s$ with $r, s$ coprime and $s \ge 1$; $m_p$ denotes the prime-to-$p$ part of an
integer $m \ne 0$, and ``$m \ne \pm\sq$'' means that neither $m$ nor $-m$ is a square.

\section{Three lists of criteria}

By Lemma~\ref{lem:klein}, $\phi_0$ and $\phi_1$ are available as soon as two of the three
patterns $\phi_0 = ({-}{-}{-}\cdots)$, $\phi_1 = ({-}{+}{+}\cdots)$, $\tau = ({+}{-}{-}\cdots)$
are. Each of the following conditions provides one of them.

\begin{proposition}[Criteria for $\phi_0$]
  \label{prop:Z}
  \uses{def:phin, lem:pat-r, lem:pat-r2s, lem:pat-r2s2}
  \lean{QuadraticIterates.NegAllCriterion,
    QuadraticIterates.hasCharPattern_neg_one_of_negAllCriterion,
    QuadraticIterates.hasCharPattern_neg_one_of_not_isSquare,
    QuadraticIterates.hasCharPattern_neg_one_of_dvd_add_two_mul,
    QuadraticIterates.hasCharPattern_neg_one_of_dvd_sq_add_sq,
    QuadraticIterates.hasCharPattern_neg_one_of_dvd_left,
    QuadraticIterates.hasCharPattern_neg_one_of_dvd_left_two}
  \leanok
  $a$ has $\phi_0$ if one of the following holds.
  \begin{enumerate}
    \item[(Z1)] $r \ne \pm\sq$.
    \item[(Z2)] A prime $p \equiv \pm 3 \pmod 8$ divides $r + 2s$.
    \item[(Z3)] An odd prime $p$ divides $r^2 + s^2$, and $\Leg{-rs}{p} = \Leg{(r - s)s}{p} = -1$.
    \item[(Z4)] A prime $p \equiv 1 \pmod 4$ divides $r$ with $\Leg{r_p s}{p} = -1$; or $8 \mid r$
          and $r_2 s \equiv \pm 3 \pmod 8$.
  \end{enumerate}
  \formalnote{An inductive proposition \lname{NegAllCriterion r s} with one constructor per item
  (two for (Z4)), stated with Mathlib's notions only: Legendre symbols as \lname{jacobiSym} and
  $p$-parts as $r = p^k m$ with $p \nmid m$.}
\end{proposition}
\begin{proof}
  \leanok
  (Z1) Some prime has $v_p(r)$ odd; $\val_p$ (Lemma~\ref{lem:pat-r}(1)). (Z2) $\chi_p$
  (Lemma~\ref{lem:pat-r2s}). (Z3) $\chi_p$ (Lemma~\ref{lem:pat-r2s2}). (Z4) $\chi_p$
  (Lemma~\ref{lem:pat-r}(2), $\Leg{-1}{p} = 1$), resp.\ $\chi_8$ (Lemma~\ref{lem:pat-r}(4)).
\end{proof}

\begin{proposition}[Criteria for $\phi_1$]
  \label{prop:O}
  \uses{lem:pat-sign, lem:pat-dyadic, lem:pat-s, lem:pat-r}
  \lean{QuadraticIterates.NegFirstCriterion,
    QuadraticIterates.hasCharPattern_mulSingle_one_of_negFirstCriterion,
    QuadraticIterates.hasCharPattern_mulSingle_one_of_dvd_right_of_odd,
    QuadraticIterates.hasCharPattern_mulSingle_one_of_dvd_right,
    QuadraticIterates.hasCharPattern_mulSingle_one_of_dvd_right_two,
    QuadraticIterates.hasCharPattern_mulSingle_one_of_dvd_right_four,
    QuadraticIterates.hasCharPattern_mulSingle_one_of_dvd_right_eight,
    QuadraticIterates.hasCharPattern_mulSingle_one_of_dvd_left,
    QuadraticIterates.hasCharPattern_mulSingle_one_of_dvd_left_four}
  \leanok
  $a$ has $\phi_1$ if one of the following holds.
  \begin{enumerate}
    \item[(O1)] $a > 0$.
    \item[(O2)] $a \equiv 1 \bmod 4$, i.e.\ $s$ is odd and $r \equiv s \pmod 4$.
    \item[(O3)] An odd prime $p$ divides $s$ with $v_p(s)$ odd or $\Leg{-r s_p}{p} = -1$.
    \item[(O4)] $v_2(s)$ is odd, or $v_2(s) = 2$ and $r s_2 \equiv 1 \pmod 4$, or $v_2(s) \ge 4$ is
          even and $r s_2 \not\equiv -1 \pmod 8$.
    \item[(O5)] A prime $p \equiv 3 \pmod 4$ divides $r$ with $\Leg{r_p s}{p} = 1$; or $4 \mid r$
          and $r_2 s \equiv 1 \pmod 4$.
  \end{enumerate}
  In particular, $a$ has $\phi_1$ if $s$ is not a square (O3 or O4).
  \formalnote{\lname{NegFirstCriterion r s}. In (O4), $\chi_4$ works whenever $4 \mid s$ and
  $r s_2 \equiv 1 \pmod 4$, and $\chi_8$ whenever $8 \mid s$ and $r s_2 \equiv \pm 3 \pmod 8$.}
\end{proposition}
\begin{proof}
  \leanok
  \uses{lem:sqclass}
  (O1) the sign (Lemma~\ref{lem:pat-sign}). (O2) Lemma~\ref{lem:pat-dyadic}. (O3) $\val_p$ or
  $\chi_p$ (Lemma~\ref{lem:pat-s}(1)). (O4) By Lemma~\ref{lem:pat-s}(2) the characters at $2$
  are trivial on $c_n$, $n \ge 3$, so it suffices to find one with value $-1$ at $c_1$ and $1$ at
  $c_2$; here $c_1 \sim -rs$, $c_2 \sim r(r+s)$ (Lemma~\ref{lem:sqclass}), $r$ odd. If $v_2(s)$ is
  odd, $\val_2$ works. If $v_2(s) = 2$, then $r + s \equiv r + 4 \pmod 8$, so
  $\chi_4(c_2) = \chi_4(r^2) = 1$, and $\chi_4(c_1) = \chi_4(-r s_2) = -1$. If $v_2(s) \ge 4$, then
  $r(r + s) \equiv 1 \pmod 8$, so every character at $2$ is $1$ at $c_2$, and as $-r s_2$ is a
  $2$-adic unit $\not\equiv 1 \pmod 8$ one of $\chi_4$, $\chi_8$, $\chi_4\chi_8$ is $-1$ at $c_1$.
  (O5) $\chi_p$, resp.\ $\chi_4$ (Lemma~\ref{lem:pat-r}(2), (3)).
\end{proof}

\begin{proposition}[Criteria for $\tau$]
  \label{prop:T}
  \uses{lem:pat-sign, lem:pat-r, lem:pat-rs}
  \lean{QuadraticIterates.NegTailCriterion,
    QuadraticIterates.hasCharPattern_neg_mulSingle_one_of_negTailCriterion,
    QuadraticIterates.hasCharPattern_neg_mulSingle_one_of_neg_lt,
    QuadraticIterates.hasCharPattern_neg_mulSingle_one_of_dvd_left,
    QuadraticIterates.hasCharPattern_neg_mulSingle_one_of_dvd_left_four,
    QuadraticIterates.hasCharPattern_neg_mulSingle_one_of_dvd_add,
    QuadraticIterates.hasCharPattern_neg_mulSingle_one_of_dvd_add_of_valChar,
    QuadraticIterates.hasCharPattern_neg_mulSingle_one_of_dvd_add_prime,
    QuadraticIterates.hasCharPattern_neg_mulSingle_one_of_dvd_add_four,
    QuadraticIterates.hasCharPattern_neg_mulSingle_one_of_not_isSquare_add,
    QuadraticIterates.hasCharPattern_neg_mulSingle_one_of_window,
    QuadraticIterates.mul_wSeq_three, QuadraticIterates.mul_wSeq_four}
  \leanok
  $a$ has $\tau$ if one of the following holds.
  \begin{enumerate}
    \item[(T1)] $-1 < a < 0$.
    \item[(T2)] A prime $p \equiv 3 \pmod 4$ divides $r$ with $\Leg{r_p s}{p} = -1$; or $4 \mid r$
          and $r_2 s \equiv 3 \pmod 4$.
    \item[(T3)] A prime $p \equiv 3 \pmod 4$ divides $r + s$ with $\Leg{s (r+s)_p}{p} = 1$; or
          $4 \mid r + s$ and $s (r+s)_2 \equiv 1 \pmod 4$; or $r + s \ne \pm\sq$ and $r + s$ is
          divisible by $4$ or by a prime $p \equiv 3 \pmod 4$.
    \item[(T4)] $a < -1$, $c_3 < 0 < c_4$ (i.e.\ $a_* < a < -1$), and $r + s \ne \pm\sq$.
  \end{enumerate}
  \formalnote{\lname{NegTailCriterion r s}; in (T4) the window $c_3 < 0 < c_4$ is stated as
  $r w_3 < 0 < r w_4$, with $r w_3 = (r^2 + rs)^2 + rs^3$ and $r w_4 = (r w_3)^2 + rs^7$.}
\end{proposition}
\begin{proof}
  \leanok
  \uses{lem:pat-rs}
  (T1) the sign (Lemma~\ref{lem:pat-sign}). (T2) $\chi_p$ resp.\ $\chi_4$
  (Lemma~\ref{lem:pat-r}(2), (3)). (T3) Lemma~\ref{lem:pat-rs}: $\chi$ itself when $y = -1$,
  and $\chi \val_q$ for a prime $q$ with $v_q(r + s)$ odd otherwise. (T4) $\sgn \cdot \val_q$ with
  $v_q(r + s)$ odd: $({+}{+}{-}{+}{-}\cdots) \cdot ({+}{-}{+}{-}\cdots) = ({+}{-}{-}{-}\cdots)$
  (Lemmas~\ref{lem:pat-sign}(3), \ref{lem:pat-rs}(1)).
\end{proof}

\section{Reflection at every level}

\begin{proposition}[Surjectivity by reflection]
  \label{prop:crit-reflection}
  \uses{lem:reflection, lem:jump, thm:tfae}
  \lean{QuadraticIterates.forall_nonempty_mulEquiv_of_forall_dvd_vSeq}
  \leanok
  Suppose that $a$ has $\phi_0$ and $\phi_1$, and that for every $n \ge 2$ the number $v_n$
  is divisible by a prime $p \equiv 3 \pmod 4$ not dividing $r$, or by $4$ with $r$ odd. Then
  $\rho_a$ is surjective.
\end{proposition}
\begin{proof}
  \leanok
  Lemma~\ref{lem:reflection} gives $\phi_n$ for $n \ge 2$; Theorem~\ref{thm:tfae}.
\end{proof}

A sufficient condition for the hypothesis on $v_n$ is $|v_n| \equiv 3 \pmod 4$ or
$|v_n| \equiv 6 \pmod 8$ (Lemma~\ref{lem:prime3}, and Lemma~\ref{lem:wseq-res} for $p \nmid r$).

\section{The \texorpdfstring{$2$}{2}-adic condition}

\begin{lemma}[Prime divisors in residue classes]
  \label{lem:prime3}
  \lean{Nat.exists_prime_dvd_of_mod_four_eq_three, Nat.exists_prime_dvd_of_mod_eight_eq_six,
    Nat.exists_prime_dvd_of_mod_eight_eq_three_or_five}
  \leanok
  Let $m > 0$. If $m \equiv 3 \pmod 4$ or $m \equiv 6 \pmod 8$, then a prime $p \equiv 3 \pmod 4$
  divides $m$. If $m \equiv \pm 3 \pmod 8$, then a prime $p \equiv \pm 3 \pmod 8$ divides $m$.
\end{lemma}
\begin{proof}
  \leanok
  The residues $\{1\}$ modulo $4$ and $\{\pm 1\}$ modulo $8$ are closed under multiplication, so
  a product of primes in them stays in them; for $m \equiv 6 \pmod 8$ apply the first statement to
  $m/2$.
\end{proof}

\begin{lemma}[The $2$-adic condition]
  \label{lem:twoadic}
  \uses{def:calA, lem:reduction, lem:wseq-res, lem:signs, lem:sumsign, lem:prime3}
  \lean{QuadraticIterates.TwoAdicClass, QuadraticIterates.exists_prime_dvd_vSeq_of_twoAdicClass,
    QuadraticIterates.mul_vSeq_pos, QuadraticIterates.vSeq_modEq_add_two_mul,
    QuadraticIterates.vSeq_modEq_two_mul_add, QuadraticIterates.vSeq_modEq_of_even}
  \leanok
  Let $a \in \cA$, $\eps = \eps_a$, and assume one of
  \begin{enumerate}
    \item[(C1)] $s$ is odd and $r + 2s \equiv -\eps \pmod 4$;
    \item[(C2)] $s$ is odd and $r + s \equiv -\eps \pmod 4$;
    \item[(C3)] $s$ is even and $r \equiv -\eps \pmod 4$.
  \end{enumerate}
  Then for every $n \ge 2$ an odd prime $p \equiv 3 \pmod 4$ divides $v_n$, and $p \nmid r$.
  \formalnote{The condition is \lname{TwoAdicClass r s ε}; the sign, $\eps v_n > 0$, is
  \lname{mul\_vSeq\_pos}.}
\end{lemma}
\begin{proof}
  \leanok
  \emph{Sign.} By Lemma~\ref{lem:reduction}, $c_n + c_{n+1} = r v_n / s^{2^n}$. For $a > 0$, $r > 0$ and
  $c_n, c_{n+1} > 0$; for $\alpha < a < 0$, $r < 0$ and $c_n + c_{n+1} < 0$
  (Lemma~\ref{lem:sumsign}); for $a \le -2$, $r < 0$ and $c_n, c_{n+1} > 0$. So $\eps v_n = |v_n| > 0$.

  \emph{Residues.} If $s$ is odd, then $s^{2^{n-1}} \equiv 1 \pmod 8$, so
  $v_n \equiv w_{n+1} + w_n \pmod 8$, and $w_{m+1} \equiv r w_m^2 + s \pmod 8$
  (Lemma~\ref{lem:wseq-res}). If moreover $r$ is odd, then $w_{2k} \equiv r + s$ and
  $w_{2k+1} \equiv s \pmod 4$ for $k \ge 1$, so $v_n \equiv r + 2s \pmod 4$ and
  $|v_n| \equiv \eps(r + 2s) \equiv 3 \pmod 4$ under (C1). If $r$ is even, then
  $w_m \equiv r + s \pmod 8$ for $m \ge 2$, so $v_n \equiv 2(r + s) \pmod 8$ and
  $|v_n| \equiv 6 \pmod 8$ under (C2). If $s$ is even, then $v_n \equiv w_{n+1} \equiv r \pmod 4$ and
  $|v_n| \equiv 3 \pmod 4$ under (C3). (In (C1) $r$ is necessarily odd, in (C2) even.)

  Lemma~\ref{lem:prime3} gives $p$, and $p \nmid r$ by Lemma~\ref{lem:wseq-res}.
\end{proof}

\section{The main theorem}

\begin{proposition}[Main theorem, abstract form]
  \label{prop:main-abstract}
  \uses{lem:twoadic, prop:crit-reflection}
  \lean{QuadraticIterates.nonempty_mulEquiv_of_twoAdicClass_of_hasCharPattern}
  \leanok
  Let $a \in \cA$ satisfy (C1), (C2) or (C3), and suppose that $a$ has $\phi_0$ and $\phi_1$.
  Then $a$ has $\phi_n$ for all $n \ge 1$, so $\rho_a$ is surjective.
\end{proposition}
\begin{proof}
  \leanok
  Lemma~\ref{lem:twoadic} supplies the primes required by Proposition~\ref{prop:crit-reflection}.
\end{proof}

\begin{theorem}[Main theorem]
  \label{thm:main}
  \uses{prop:main-abstract, prop:Z, prop:O, prop:T, lem:klein}
  \lean{QuadraticIterates.nonempty_mulEquiv_of_twoAdicClass_of_criteria}
  \leanok
  Let $a = r/s \in \QQ$ with $r, s$ coprime and $s \ge 1$, and let either $a \le -2$ and $\eps = -1$,
  or $a^3 + 2a^2 + 2a + 2 > 0$ and $\eps = 1$. Assume that one of (C1), (C2),
  (C3) of Lemma~\ref{lem:twoadic} holds, and that at least two of the following hold: one of the
  conditions (Z1)--(Z4) of Proposition~\ref{prop:Z}; one of the conditions (O1)--(O5) of
  Proposition~\ref{prop:O}; one of the conditions (T1)--(T4) of Proposition~\ref{prop:T}. Then
  $\Om_n \cong \Ct{n}$ for all $n$: the arboreal representation $\rho_a$ is surjective.
  \formalnote{Stated with the definitions \lname{NegAllCriterion}, \lname{NegFirstCriterion},
  \lname{NegTailCriterion} and \lname{TwoAdicClass}, which use Mathlib's notions only. The
  exclusion of $a = -1, 0$ from $\cA$ is not needed: $\phi_0$ forces $c_1 = -a \ne 0$ and
  $c_2 = a(a + 1) \ne 0$.}
\end{theorem}
\begin{proof}
  \leanok
  Propositions~\ref{prop:Z}, \ref{prop:O}, \ref{prop:T} and Lemma~\ref{lem:klein} give $\phi_0$
  and $\phi_1$; apply Proposition~\ref{prop:main-abstract}.
\end{proof}

\begin{remark}[Numerical evidence]
  In the box $|r| \le 250$, $1 \le s \le 50$, Theorem~\ref{thm:main} covers $7399$ parameters, and
  none of them has a square product among $c_1, \dots, c_6$. Each item of the three lists was
  checked to produce the claimed pattern on $c_1, \dots, c_8$ for all $5518$ parameters with
  $|r| \le 150$, $s \le 30$.
\end{remark}

\begin{corollary}[Li 2020 for rational parameters]
  \label{cor:li2020-rat}
  \uses{prop:main-abstract, prop:O, prop:T, lem:klein, thm:tfae}
  \lean{QuadraticIterates.forall_nonempty_mulEquiv_iff_of_le_neg_two}
  \leanok
  Let $a \le -2$ with $s \equiv 1 \pmod 4$ and $r \equiv 3 \pmod 4$, and suppose that $a$ has
  $\phi_0$ (e.g.\ $|r|$ is not a square). Then $\rho_a$ is surjective if and only if $-(a + 1)$
  is not a square.
\end{corollary}
\begin{proof}
  \leanok
  If $-(a + 1)$ is a square, so is $c_1 c_2 = -a \cdot a(a + 1)$, and $\rho_a$ is not surjective
  (Theorem~\ref{thm:tfae}). Otherwise: (C1) holds, as $r + 2s \equiv 1 = -\eps \pmod 4$. If $s$ is
  not a square, (O3) gives $\phi_1$. If $s$ is a square, then $4 \mid r + s = s(a + 1) < 0$ and
  $-(r + s) = s \cdot (-(a+1))$ is not a square, so (T3) gives $\tau$, and $\phi_1$ follows by
  Lemma~\ref{lem:klein}. Apply Proposition~\ref{prop:main-abstract}.
\end{proof}

\section{When the second level decides}

The classes (C1)--(C3) of Lemma~\ref{lem:twoadic} are, in terms of $(r, s) \bmod 4$: for
$\eps = -1$, $(3,1)$, $(3,3)$; $(0,1)$, $(2,3)$; $(1,0)$, $(1,2)$; and for $\eps = 1$, $(1,1)$,
$(1,3)$; $(2,1)$, $(0,3)$; $(3,0)$, $(3,2)$. Outside the interval $(\alpha, a_*]$, the three
lists of criteria collapse to a single condition.

\begin{corollary}[Level two decides]
  \label{cor:level2}
  \uses{thm:main, thm:criterion, prop:Z, prop:O, prop:T, lem:prime3, lem:window}
  \lean{QuadraticIterates.forall_nonempty_mulEquiv_iff_twoIndependent_of_twoAdicClass,
    QuadraticIterates.not_isSquare_of_twoIndependent, QuadraticIterates.twoIndependent_fin_two_iff,
    QuadraticIterates.isSquare_cSeq_one_of_isSquare, QuadraticIterates.isSquare_cSeq_two_of_isSquare,
    QuadraticIterates.isSquare_cSeq_one_mul_cSeq_two_of_isSquare, Int.sq_emod_eight_eq_one_of_odd}
  \leanok
  Let $a = r/s$ with $r, s$ coprime and $s \ge 1$, and let either $a \le -2$ and $\eps = -1$, or
  $a > a_*$, $a \ne -1, 0$, and $\eps = 1$, with $a_*$ as in Lemma~\ref{lem:window}. Assume
  that one of (C1), (C2), (C3) holds. Then $\rho_a$ is surjective if and only if $c_1, c_2$ are
  $2$-independent, that is, $\Om_2 \cong \Ct{2}$, that is, none of $-a$, $a(a+1)$ and
  $-a^2(a+1)$ is a square.
  \formalnote{The condition is \lname{TwoIndependent (fun i : Fin 2 ↦ cSeq a (i + 1))}, and
  $a > a_*$ is stated as $-s < r$, or $r < -s$ and $r w_3 < 0 < r w_4$ as in (T4). That this
  window lies above $\alpha$ is shown with $X = r^2 + rs$ and $Y = X^2 + rs^3$: $Y < -s^4$, hence
  $X < s^2$, and $r(r^3 + 2r^2 s + 2rs^2 + 2s^3) = X s^2 + Y < 0$.}
\end{corollary}
\begin{proof}
  \leanok
  \uses{thm:criterion, thm:main, prop:Z, prop:O, prop:T, lem:prime3, lem:window}
  The condition is necessary by Theorem~\ref{thm:criterion}. Modulo squares, $c_1 \sim -rs$,
  $c_2 \sim r(r+s)$ and $c_1 c_2 \sim -s(r+s)$, and $r$, $s$, $r + s$ are pairwise coprime. So
  $c_1$ is a square if and only if $r < 0$ and $|r|$, $s$ are squares; $c_2$ if and only if $r$,
  $r + s$ have the same sign and $|r|$, $|r + s|$ are squares; $c_1 c_2$ if and only if $r + s < 0$
  and $s$, $|r + s|$ are squares. We show that two of (Z), (O), (T) hold; as $a_* > \alpha$,
  Theorem~\ref{thm:main} then applies. Recall that (Z1) holds if $|r|$ is not a square, (O3) or
  (O4) if $s$ is not a square, and (T3) if $|r + s|$ is not a square and $|r + s| \equiv 0$ or
  $3 \pmod 4$ (Lemma~\ref{lem:prime3}).

  \emph{$-1 < a < 0$.} (T1) holds, and $|r|$ or $s$ is not a square: (Z1) or (O3)/(O4).

  \emph{$a_* < a < -1$.} At most one of $|r|$, $s$, $|r + s|$ is a square; the other two give
  two of (Z1), (O3)/(O4) and (T4).

  \emph{$a > 0$.} (O1) holds, and $r$ or $r + s$ is not a square. In the classes $(1,3)$,
  $(2,1)$, $(0,3)$, $(3,0)$, $r + s \equiv 0$ or $3 \pmod 4$, so this gives (Z1) or (T3). In the
  class $(3,2)$, $r \equiv 3 \pmod 4$ is not a square: (Z1). In the class $(1,1)$, (Z1) holds
  unless $r = m^2$; then $r \equiv 1 \pmod 8$ and $r + 2s \equiv 3 \pmod 8$, and (Z2) holds.

  \emph{$a \le -2$.} At most one of $|r|$, $s$, $|r + s|$ is a square. In the classes $(3,1)$,
  $(0,1)$, $(2,3)$, $(1,0)$, $|r + s| \equiv 0$ or $3 \pmod 4$, and the two non-squares give two
  of (Z1), (O3)/(O4) and (T3). In the class $(3,3)$, (O2) holds, and (Z1) unless $r = -m^2$; then
  $r \equiv 7$ and $r + 2s \equiv 5 \pmod 8$, and (Z2) holds. In the class $(1,2)$,
  $|r| \equiv 3$ and $s \equiv 2 \pmod 4$ are not squares: (Z1) and (O4).
\end{proof}

\begin{remark}[When the condition is automatic]
  \label{rem:level2-classes}
  By the residues of $|r|$, $s$ and $|r + s|$ modulo $4$, the condition of
  Corollary~\ref{cor:level2} holds automatically, and $\rho_a$ is surjective, except in the
  following cases, where it says:
  \begin{itemize}
    \item $a \le -2$: in the class $(0,1)$, $-a \ne \sq$; in the class $(3,1)$, at most one of
          $|r|$, $s$, $|r + s|$ is a square (a sharper form of Corollary~\ref{cor:li2020-rat});
    \item $a > 0$: in the class $(1,3)$, $a(a+1) \ne \sq$;
    \item $-1 < a < 0$: in the class $(3,0)$, $-a \ne \sq$;
    \item $a_* < a < -1$: in the class $(2,1)$, $-(a+1) \ne \sq$; in the class $(0,3)$,
          $a(a+1) \ne \sq$; in the class $(3,0)$, at most one of $|r|$, $s$, $|r + s|$ is a square.
  \end{itemize}
  For $s = 1$ the classes are $a \equiv 1, 2 \pmod 4$ with $a > 0$ and $a \equiv 0, 3 \pmod 4$
  with $a \le -2$, and the corollary contains Corollaries~\ref{cor:stoll1992}
  and~\ref{cor:li2020}: $\rho_a$ is surjective if and only if neither $-a$ nor $-a - 1$ is a
  square. Numerically, for all $469810$ parameters with $|r| \le 2000$, $s \le 400$ in these
  classes and outside $(\alpha, a_*]$, the hypotheses of Theorem~\ref{thm:main} hold exactly when
  $c_1, c_2$ are $2$-independent.
\end{remark}

\begin{remark}[Square products among $c_1, \dots, c_{10}$]
  \label{rem:square-products}
  As polynomials in $a$, $c_1 = -b_1$ and $c_n = \prod_{d \mid n} b_d$ for $n \ge 2$, where
  $b_1 = a$, $b_2 = a + 1$, $b_3 = a^3 + 2a^2 + a + 1$, $b_4 = a^6 + 3a^5 + 3a^4 + 3a^3 + 2a^2 + 1$,
  and in general the roots of $b_d$ are the $a$ for which $0$ is periodic of exact period $d$. For
  $d \le 10$ the $b_d$ are irreducible, of degrees $1, 1, 3, 6, 15, 27, 63, 120, 252, 495$. For a
  nonempty $S \subseteq \{1, \dots, 10\}$ let $D$ be the set of $d$ that divide an odd number of
  elements of $S$, and $\eps_S = -1$ if $1 \in S$, $\eps_S = 1$ otherwise. For $a \ne -1, 0$, the
  product $\prod_{n \in S} c_n$ is a square if and only if $\eps_S \prod_{d \in D} b_d(a)$ is, so
  these $a$ are the abscissae of the rational points on $y^2 = \eps_S \prod_{d \in D} b_d(a)$. As
  $\max D = \max S$, this curve has genus $0$ exactly for $S \subseteq \{1, 2\}$, and genus $1$
  exactly for $S = \{3\}$, $\{1, 3\}$, $\{2, 3\}$. A search with \texttt{ratpoints}
  \cite{ratpoints}, its degree bound raised to $1024$, for the points with $a = r/s$,
  $\max(|r|, s) \le 10^5$, on the $1020$ curves of positive genus finds some with
  $a \notin \{-2, -1, 0\}$ in three cases only. The nonempty lists are:
  \begin{itemize}
    \item $S = \{1\}$: $a = -t^2$; $S = \{2\}$: $a = 1/(t^2 - 1)$; $S = \{1, 2\}$: $a = -1 - t^2$,
          with $t \in \QQ$;
    \item $S = \{1, 3\}$, $y^2 = -b_3(a)$, the elliptic curve 368d1 of rank $1$: $a = -17/4$,
          $-586/81$, $-81185/44944$;
    \item $S = \{2, 3\}$, $y^2 = (a + 1)\,b_3(a)$, the elliptic curve 92b1 of rank $1$: $a = 3$,
          $-2/3$, $-4/5$, $6/19$, $-20/11$, $-55/56$, $-38/159$, $-616/255$, $-1007/223$,
          $-2856/5665$, $-23638/26239$, $26265/23464$;
    \item $S = \{2, 4\}$, $y^2 = b_4(a)$, of genus $2$, whose Jacobian has rank $3$ (by
          $2$-descent), so that Chabauty's method does not apply: $a = 2/3$, $-6/7$.
  \end{itemize}
  For $S = \{3\}$, $y^2 = a\,b_3(a)$ is the elliptic curve 92a1 of rank $0$, and $c_3$ is a square
  only for $a = 0$. So for $a \notin \{-2, -1, 0\}$ with $\max(|r|, s) \le 10^5$, the numbers
  $c_1, \dots, c_{10}$ are $2$-dependent if and only if one of $-a$, $a(a + 1)$, $-(a + 1)$ is a
  square or $a$ is one of the $17$ values listed; for all other such $a$,
  $\Om_{10} \cong \Ct{10}$ (Theorem~\ref{thm:criterion}).
\end{remark}

\section{\texorpdfstring{The interval $(\alpha, -1)$}{The interval (alpha, -1)}}

\begin{corollary}[The interval]
  \label{cor:interval}
  \uses{thm:main, prop:Z, prop:O, prop:T, cor:level2}
  \lean{QuadraticIterates.nonempty_mulEquiv_of_lt_neg_one_of_criteria,
    QuadraticIterates.NegFirstCriterion', QuadraticIterates.NegTailCriterion',
    QuadraticIterates.negFirstCriterion_of_negFirstCriterion',
    QuadraticIterates.negTailCriterion_of_negTailCriterion', jacobiSym.sq_one_of_not_dvd,
    Int.exists_sq_eq_pow_mul_of_dvd}
  \leanok
  Let $a = r/s$ with $r, s$ coprime, $s \ge 1$ and $\alpha < a < -1$, assume that one of (C1),
  (C2), (C3) holds with $\eps = 1$, and that $c_1, c_2$ are $2$-independent. Consider the
  conditions
  \begin{enumerate}
    \item[(O$'$)] $a \equiv 1 \bmod 4$; or $s \ne \sq$; or $\Leg{-r}{p} = -1$ for an odd prime
          $p \mid s$; or $16 \mid s$ and $r \equiv 3 \pmod 8$; or a prime $p \equiv 3 \pmod 4$
          divides $r$ with $\Leg{r_p s}{p} = 1$;
    \item[(T$'$)] a prime $p \equiv 3 \pmod 4$ divides $r$ with $\Leg{r_p s}{p} = -1$; or a prime
          $p \equiv 3 \pmod 4$ divides $r + s$.
  \end{enumerate}
  If two of (Z), (O$'$), (T$'$) hold, where (Z) is one of (Z1)--(Z4) of
  Proposition~\ref{prop:Z}, then $\rho_a$ is surjective.
  \formalnote{(O$'$) and (T$'$) are the inductive propositions \lname{NegFirstCriterion'} and
  \lname{NegTailCriterion'}; the two implications of the proof are
  \lname{negFirstCriterion\_of\_negFirstCriterion'} and
  \lname{negTailCriterion\_of\_negTailCriterion'}, the latter for $|r + s| \ne \sq$.}
\end{corollary}
\begin{proof}
  \leanok
  \uses{thm:main, prop:Z, prop:O, prop:T}
  (O$'$) implies (O): (O2); (O3) or (O4) if $s \ne \sq$; if $s$ is a square, then so is $s_p$,
  and $\Leg{-r s_p}{p} = \Leg{-r}{p}$ gives (O3), and $16 \mid s = 2^k s_2$ forces $k \ge 4$ even
  and $s_2 \equiv 1 \pmod 8$, so $r s_2 \equiv 3 \pmod 8$ gives (O4); and (O5). (T$'$) implies
  (T) when $|r + s|$ is not a square: (T2), and (T3), as $r + s < 0$. If $|r + s|$ is a square,
  then neither $|r|$ nor $s$ is ($c_2 \sim |r|\,|r + s|$ and $c_1 c_2 \sim s\,|r + s|$ are not
  squares), so (Z1) and (O3)/(O4) hold. In both cases Theorem~\ref{thm:main} applies.
\end{proof}

\begin{remark}[The six classes on the interval]
  \label{rem:interval}
  Let $\alpha < a < -1$. In the classes $(1,1)$, $(1,3)$, $(3,2)$, the numbers $c_1, c_2$ are
  $2$-independent and two of (Z), (O$'$), (T$'$) hold automatically: (Z1) and $a \equiv 1 \bmod 4$;
  (Z1) and $s \equiv 3 \pmod 4$; $s \equiv 2 \pmod 4$ and $|r + s| \equiv 3 \pmod 4$. In the class
  $(2,1)$, (Z1) holds, and the condition says: $s \ne \sq$, or $\Leg{-r}{p} = -1$ for a prime
  $p \mid s$, or a prime $p \equiv 3 \pmod 4$ divides $r(r + s)$. In the class $(0,3)$, (O$'$)
  holds, and the condition says (Z) or (T$'$). In the class $(3,0)$, none of the three holds
  automatically.

  For $\alpha < a < a_*$ and $c_1, c_2$ $2$-independent, two of (Z), (O$'$), (T$'$) hold if and
  only if two of (Z), (O), (T) do: there (O1), (T1) and (T4) do not apply, (O3), (O4) and (O5)
  amount to (O$'$), and the parts of (T2) and (T3) at $2$ either do not apply or occur only
  together with (Z1). Numerically, the two agree for all $318538$ such parameters with
  $s \le 3000$. Unlike outside the interval, the criteria are not sharp here: for $s \le 2500$,
  $127$ parameters in the classes $(2,1)$, $(0,3)$, $(3,0)$ fail them although
  $c_1, \dots, c_6$ are $2$-independent, for instance $a = -634/441$, $-225/164$ and $-1296/887$.
\end{remark}

\begin{remark}[Density]
  \label{rem:density}
  Order the parameters $a = r/s$ by the height $\max(|r|, s)$. Asymptotically, the coprime pairs
  $(r, s)$ are equidistributed among the $12$ residue classes modulo $4$ in which $r$ or $s$ is
  odd, those with $|r|$ or $s$ a square have density zero, and $a$ is distributed with density
  $\frac14$ on $[-1, 1]$ and $\frac{1}{4a^2}$ outside it. Theorem~\ref{thm:main} applies whenever
  $a \in \cA$ lies in one of the six classes (C1)--(C3) for $\eps_a$ and neither $|r|$ nor $s$ is a
  square: then (Z1) holds, and (O3) or (O4). So it covers the proportion $\frac12$ of the
  parameters in every interval contained in $\cA$, in $(\alpha, a_*]$ just as outside it (the
  interval differs only on the density-zero set where the criteria are not sharp), and none in
  $(-2, \alpha]$, which carries the proportion $\frac{1}{4|\alpha|} - \frac18 \approx 0.0370$.
  Altogether the main theorem covers the proportion $\frac{9}{16} - \frac{1}{8|\alpha|} \approx
  0.4815$ of all $a \in \QQ$; the families of Chapter~\ref{chap:families}, with $s$ a square, have
  density zero. Numerically, for $\max(|r|, s) \le 400$ the two proportions are $0.0369$ and
  $0.4796$.
\end{remark}

\section{Odd levels for \texorpdfstring{$a \equiv 1 \bmod 4$}{a = 1 mod 4}}

\begin{corollary}[Odd levels]
  \label{cor:odd}
  \uses{lem:period2, lem:pat-dyadic, lem:jump2, lem:reduction, lem:signs, lem:window, lem:prime3}
  \lean{QuadraticIterates.hasCharPattern_mulSingle_of_odd,
    QuadraticIterates.exists_prime_dvd_uSeq_of_odd, QuadraticIterates.uSeq_emod_eight_of_odd,
    QuadraticIterates.uSeq_two_mul_add_one_modEq, QuadraticIterates.uSeq_pos,
    QuadraticIterates.cSeq_sq_sub_cSeq_add_pos_of_le_neg_two,
    QuadraticIterates.cSeq_sq_sub_cSeq_add_pos_of_neg_one_lt,
    QuadraticIterates.cSeq_sq_sub_cSeq_add_pos_of_odd}
  \leanok
  Let $a \equiv 1 \bmod 4$ with $a \le -2$, or $a > -1$, or $a < -1$ and $c_3 < 0 < c_4$. If $a$
  has $\phi_0$, then $a$ has $\phi_n$ for every odd $n \ge 3$.
  \formalnote{For $a = r/s$, the condition $a \equiv 1 \bmod 4$ is stated as $s$ odd and
  $r \equiv s \pmod 4$.}
\end{corollary}
\begin{proof}
  \uses{lem:pat-rs}
  \leanok
  Let $n \ge 3$ be odd. In $\ZZ/8\ZZ$, $\bar c_n = 5$ for $\bar a = 1$ and $\bar c_n = 1$ for
  $\bar a = 5$, and $s^{2^n} \equiv 1$, so $u_n \equiv \bar c_n^2 - \bar c_n + \bar a + 1 \equiv 6
  \pmod 8$ (Lemma~\ref{lem:reduction}). Also $u_n > 0$: $u_n / s^{2^n} = c_n^2 - c_n + a + 1$, which is
  positive for $a > 0$ (no real root), for $-1 < a < 0$ (as $c_n < 0$), in the window (as
  $c_n \le c_3 < -1$, Lemma~\ref{lem:window}), and for $a \le -2$ (as $c_n \ge -a \ge 2$ gives
  $c_n(c_n - 1) + a + 1 \ge (a + 1)^2$). So a prime $p \equiv 3 \pmod 4$ divides $u_n$
  (Lemma~\ref{lem:prime3}). It does not divide $r + s$: modulo an odd prime $q \mid r + s$,
  $\bar a = -1$ and $(\bar c_m) = (1, 0, -1, 0, -1, \dots)$, so
  $\bar c_n^2 - \bar c_n + \bar a + 1 = 2 \ne 0$. Lemma~\ref{lem:period2}, with the alternating
  character $\val_2$ (Lemma~\ref{lem:pat-dyadic}) and $\phi_0$, gives $\phi_n$.
\end{proof}

\begin{remark}[The even levels]
  \label{rem:even}
  For $a \le -2$, $a \equiv 1 \bmod 4$ and $s \equiv 1 \pmod 4$ (for $s \equiv 3 \pmod 4$ the
  main theorem applies), one has $\phi_0$ (when given), $\phi_1$ (O2), $\phi_2$
  (Remark~\ref{rem:phi2}) and $\phi_n$ for odd $n$, but no argument of this kind reaches the
  even $n \ge 4$: $|v_n| \equiv 1 \pmod 4$, and $u_n$ is $4$ times an odd number whose part prime
  to $r + s$ is $\equiv 1 \pmod 4$ (checked for all integers $-1203 \le a \le -3$,
  $a \equiv 1 \pmod 4$, and $n = 4, 6, 8$, in every class modulo $64$). This is the counterpart of
  the barrier $n_0 \in \{4, 8, 16, \dots\}$ of the approach via $b_n$. The families of
  Chapter~\ref{chap:families} get past it because $x^2 + x + a$ or $x^2 - x + a + 1$ splits.
\end{remark}
